\documentclass[11pt,letterpaper]{article}
\usepackage[margin=.9in,headheight=15pt]{geometry}
\usepackage{amsmath,amssymb,amsthm,mathtools,booktabs,longtable,array,microtype,xcolor,fancyhdr,enumitem,hyperref}
\microtypesetup{expansion=false}
\definecolor{ink}{HTML}{17384A}
\definecolor{accent}{HTML}{176B65}
\hypersetup{colorlinks=true,linkcolor=accent,citecolor=accent,urlcolor=accent,pdftitle={Beyond Dirichlet GRH: orthogonal kernels and sharp p-adic recovery},pdfauthor={Research analysis prepared for Vladimir Reshetnikov}}
\pagestyle{fancy}\fancyhf{}\fancyhead[L]{\small\sffamily\color{ink}Method transfers from the GRH formalization}
\fancyhead[R]{\small\sffamily 26 November 2026}\fancyfoot[C]{\thepage}
\setlength{\parindent}{0pt}\setlength{\parskip}{5pt plus 1pt minus 1pt}
\setlength{\emergencystretch}{2em}
\setlist{nosep,leftmargin=*}
\newtheorem{theorem}{Theorem}[section]
\newtheorem{lemma}[theorem]{Lemma}
\newtheorem{proposition}[theorem]{Proposition}
\newtheorem{corollary}[theorem]{Corollary}
\theoremstyle{definition}\newtheorem{definition}[theorem]{Definition}
\theoremstyle{remark}\newtheorem{remark}[theorem]{Remark}
\newcommand{\F}{\mathbb F}\newcommand{\Z}{\mathbb Z}\newcommand{\Q}{\mathbb Q}
\newcommand{\R}{\mathbb R}\newcommand{\C}{\mathbb C}\newcommand{\OO}{\mathcal O}
\newcommand{\Sp}{\operatorname{Sp}}\newcommand{\GL}{\operatorname{GL}}
\newcommand{\SO}{\operatorname{SO}}\newcommand{\rank}{\operatorname{rank}}
\newcommand{\Span}{\operatorname{span}}\newcommand{\TV}{d_{\mathrm{TV}}}
\newcommand{\Prob}{\mathbb P}\newcommand{\E}{\mathbb E}
\newcommand{\src}[2]{\href{https://github.com/lsernau/grh-lean/blob/72fead102f663d5c7a407964a13dc0cf9d977207/grh/#1}{#2}}
\newcommand{\prior}[2]{\href{https://github.com/VladimirReshetnikov/ProveIt/blob/9816644452bb42827167d6110e6d0322315a02b9/docs/incoming/#1}{#2}}
\newcommand{\code}[1]{\texttt{\small #1}}
\newcolumntype{P}[1]{>{\raggedright\arraybackslash}p{#1}}
\begin{document}
\thispagestyle{plain}
{\large\sffamily\color{accent}Research report \hfill 26 November 2026}\par
\vspace{.7em}
{\fontsize{25}{29}\selectfont\bfseries\color{ink}Beyond Dirichlet GRH\par}
\vspace{.35em}
{\Large Orthogonal kernels and sharp $p$-adic recovery\par}
\vspace{.55em}
{\large Source analysis, complete mathematical proofs, and exact checks}\par
\vspace{.7em}\hrule\vspace{.7em}

\textbf{Answer to the research question.}
The repository contains useful methods beyond its final RH and Dirichlet-GRH theorems. This report develops three such transfers. The most direct answer to a separately posed question concerns random binary orthogonal matrices. If $Q$ is uniform in
\[
 \mathcal O_n=\{Q\in\GL(n,\F_2):Q^{\mathsf T}Q=I\},\qquad
 \nu(Q)=\max_{J\subseteq\{1,\ldots,n\}}\dim_{\F_2}\ker Q[J],
\]
then
\begin{equation}\label{eq:headline-binary}
 \boxed{\quad\Prob\{\nu(Q)\ge r\}\le
 \min\{1,\;6\,2^{\,n-r(r-1)/2}\}.\quad}
\end{equation}
Here $1\le r\le\lfloor n/2\rfloor$; for larger $r$ the probability is zero. Thus $\nu(Q)=O(\sqrt n)$ with exponentially high probability. In particular, for every fixed $c>0$, a linear lower bound $\nu(Q)\ge cn$ has probability tending to zero at a quadratic exponential rate, contrary to the proposed high-probability behavior in the public question of Vaas~\cite{vaas}. This settles that question's proposed linear alternative and supplies an explicit upper tail. It does not establish the exact typical order.

A second theorem answers the multibranch sharpness question explicitly left in an earlier ProveIt report~\cite{priorjets}. At precision modulo $\pi^r$, write $a_c=v_\pi(b-c)$ for the contact depths of a fixed Lubin--Tate branch, $W=\sum_{c\ne b}q^{a_c}$, $A=\max a_c$, and $m=\#\{c:a_c=A\}$. In the coordinate $[\pi](X)=\pi X+X^q$, the inverse interpolation denominator has pole order
\[
 t_b\le W+(r-1)(q-1)q^{A-1}\qquad(A\ge1).
\]
Equality holds exactly when adding $m-1$ and $r-1$ in base $p$, the residue characteristic, has no carry. We also give an exact formula at every precision when all contacts have depths zero or one, and prove the branch recovery cutoff is precisely $1+t_b$.

The third result determines an exact total-variation distance for half-size blocks of uniformly random split special orthogonal matrices over fields of odd order. A rank-parity constraint survives at that boundary; the distance from an independent-entry matrix is exactly its odd-rank probability.

\textbf{Status.} These are ordinary mathematical proofs with independently executable finite checks. They do not use RH or GRH. They are not newly kernel-checked Lean theorems. The source proof and the earlier archives are pinned below. We make no worldwide priority claim for the new formulations, and no claim to have proved nonabelian GRH, Artin holomorphy, or BSD.

\newpage
\tableofcontents
\newpage
\section{What the repository proves and what can transfer}
\label{sec:source}

\subsection{Pinned evidence and scope}

The inspected upstream commit is
\begin{center}\code{72fead102f663d5c7a407964a13dc0cf9d977207}.\end{center}
The ProveIt intake snapshot is
\begin{center}\code{9816644452bb42827167d6110e6d0322315a02b9}.\end{center}
We inventoried all 52 incoming GRH ZIP archives at that snapshot, verified their Git blob hashes, and extracted 1,630 text/code members for comparison. The directory inventory, archive hashes, and coverage notes accompany this report. Extraction and comparison are not independent validation of every prior theorem.

The upstream \src{src/Statements.lean}{statement file} defines the zero statements using Mathlib's zeta function and \code{DirichletCharacter.LFunction}. The five final exports in \src{src/Main.lean}{Main.lean} cover ordinary RH, its all-nontrivial-zero formulation, Dirichlet nonvanishing to the right of $1/2$ with the principal pole exception, and the all-character and primitive-character critical-strip formulations. These endpoints concern Dirichlet characters over $\Q$. Auxiliary CM fields and higher-rank groups inside their proofs do not enlarge the quantifiers of those conclusions.\cite{repo,statements}

The upstream README reports Lean 4.35.0 and Mathlib revision
\begin{center}\code{56a005855ea17e0b0958f73ec801ae140b0a1679},\end{center}
with dependencies only on
\code{propext}, \code{Classical.choice}, and \code{Quot.sound}, together with additional machine checks. We inspected source and declarations; we did not repeat the full Lean build. Taking the user's premise as the starting point is compatible with keeping that validation distinction explicit.\cite{repo}

\subsection{The arithmetic mechanism}

The central construction has three tightly connected estimates. Integral theta derivatives encode divisibility by prime powers. This makes normalized source forms small. A second-moment argument then ensures that enough independent target forms have been obtained. Finally, determinant identities, congruence control, and carefully selected twists prevent their corrected covolume from being too small.

\begin{center}\small
\begin{tabular}{P{.18\linewidth}P{.38\linewidth}P{.34\linewidth}}
\toprule
Stage & Precise role & Source landmark\\
\midrule
Short forms & Prime-sensitive divisibility and an upper norm/index bound & \code{ShortAndEnough.lean}, Theorem A\\[4pt]
Enough forms & Independent vectors occupying at least $1-4/b^2$ of the target dimension & \code{Data.enough\_forms}, Theorem B\\[4pt]
Long integral forms & A lower bound for the same covolume after paying for index and twist heights & \code{integral\_forms\_long}, Theorem C\\[4pt]
Cancellation & Matching large main terms leave a bound for the desired prime remainder & \code{star\_inequality}\\[4pt]
Zero exclusion & An auxiliary twist removes cancellation; a one-sided transform argument excludes poles & \code{Zeros/Criterion.lean}\\
\bottomrule
\end{tabular}
\end{center}

The explicit comparison is
\[
 m\,\mathrm{totalRemainder}\le f a(\log b)^5.
\]
Here the number of lifts satisfies $\sqrt x\le m\le2\sqrt x$. Scale choices and the remaining estimates yield
\[
 R_{d,q}(x)\le(64\varphi(q)+3)x^{-1/2+6\epsilon}
\]
eventually, for fixed admissible $q,d$ and every sufficiently small $\epsilon>0$. The remainder centers the sum of
$\Lambda(n)(1-\chi_d(n))(1/n-1/x)$ over $n\le x$ in the class $1\bmod q$.
The residue calculation sees differences of zero multiplicities. The extra quadratic twist is selected so its contribution vanishes in the relevant region, leaving nonnegative multiplicities. This noncancellation step matters as much as the final transform.\cite{guide,mainineq,criterion}

The same outline does not prove RH for a new Euler product merely by replacing a Dirichlet character in a formula. Its divisibility interpretation, spanning estimate, determinant main term, denominator costs, and noncancellation device must all survive the change.

\subsection{The two ingredients used in this report}

The source's
\src{src/New/Spanning/SecondMoment/PositiveType.lean}{positive-type matrix lemma}
\code{New.card\_mul\_sum\_det\_eq\_zero\_le} bounds the mass of singular compressions
$A_1^{\mathsf T}\pi(g)A_0$ in a finite-group representation over a finite field. It counts kernel vectors by Fourier orthogonality and controls the resulting coefficients through stabilizer masses and Cauchy--Schwarz. The incoming corpus already extends this to rectangular matrices, frame moments, and symplectic boundary blocks.\cite{positivetype,priorrank,priorsymplectic}

Sections~\ref{sec:binary} and \ref{sec:orthogonal} continue this direction. Their proofs use direct frame-orbit counts suited to orthogonal groups. The binary result handles the characteristic-two canonical vector explicitly. The odd-characteristic boundary result retains the component and parity information that a generic stabilizer bound can lose. These are developments prompted by the source ingredient, not applications of a GRH consequence.

The source's
\src{src/New/Derivatives/Sharpness/OneCurve.lean}{one-curve sharpness lemma}
\code{New.exists\_monomial\_max\_norm\_coeff} reduces an extremal integral Taylor coefficient to a finite monomial choice. The previous branch-interpolation report transferred that finite-jet viewpoint to Lubin--Tate substitutions and left a precise cancellation question at finite precision. Section~\ref{sec:jets} answers that question through interpolation denominators and their leading Newton face.\cite{onecurve,priorjets}

\subsection{What the earlier corpus already contains}

The following are comparison landmarks, not claims independently certified here.

\begin{center}\small
\begin{tabular}{P{.29\linewidth}P{.59\linewidth}}
\toprule
Prior packet & Content relevant to this analysis\\
\midrule
\path{grh_beyond_dirichlet_2026-11-25.zip} & Rectangular factorial kernel moments and rank tails; therefore a general frame-moment extension is not new to this corpus.\\[4pt]
\path{grh_sharp_torsion_and_symplectic_blocks.zip} & Exact torsion refinements and symplectic compression laws, including boundary behavior.\\[4pt]
\path{grh_transfers_cycle_signings_and_lubin_tate.zip} & Sharp residue-field branch precision; the inverse-pole characterization modulo $\pi^r$; an explicit question about multibranch cancellations.\\[4pt]
\path{grh_arithmetic_transfers_2026-11-26.zip} & Integral congruence/intertwiner depth and related arithmetic transfers.\\[4pt]
\path{grh_ideal_densities_and_transfer_conditions_2026-11-26.zip} & Density problems and explicit limitations of noncyclotomic/Selmer transfers.\\[4pt]
Other nonabelian and symmetric-power packets & Numerous exact detecting weights and one-sided analytic criteria; the required new global arithmetic inequalities remain separate tasks.\\
\bottomrule
\end{tabular}
\end{center}

We therefore do not present another equivalent nonabelian zero criterion as a newly solved conjecture. The advances developed here have complete finite-group or local-algebra proofs and an explicit comparison with the immediately preceding results.

\section{A tail bound for every principal binary orthogonal kernel}
\label{sec:binary}

\subsection{The separately posed question}

Vaas asks about
\[
 \nu(Q)=\max_{J\subseteq[n]}\dim\ker Q[J],\qquad
 Q\text{ uniform in }\{Q\in\GL(n,\F_2):Q^{\mathsf T}Q=I\},
\]
including whether constants $c,\alpha>0$ can make
$\Prob(\nu(Q)\ge cn)\ge1-2^{-\alpha n^2}$. The retrieved public page has the question and no displayed answer. We solve the proposed linear alternative negatively and give an upper bound on the scale. We do not claim that this proves the exact typical order or establishes publication priority over all random-matrix literature.\cite{vaas}

Everything below uses the symmetric bilinear dot product over $\F_2$. This is important: replacing it silently with an orthogonal group attached to a quadratic form would change the problem.

\subsection{A quantitative symplectic incidence lemma}

\begin{lemma}[Subspaces mapped into a prescribed subspace]\label{lem:symp-incidence}
Let $V$ be a nondegenerate symplectic space of even dimension $N$ over $\F_2$. Fix subspaces $U,T\subseteq V$ of dimensions $d,t$. Let $S$ be uniform in $\Sp(V)$. For $1\le s\le N/2$,
\[
 \Prob\{\exists L\subseteq U:\dim L=s,\;S L\subseteq T\}
 \le 6\,2^{-s(N-d-t)-s(s+1)/2}.
\]
If $s>d$ or $s>t$, the event is empty. For $s=0$ the bound with right side $6$ also holds.
\end{lemma}

\begin{proof}
Fix an independent ordered $s$-frame $x_1,\ldots,x_s$ spanning one candidate $L$. Its orbit consists of all independent $s$-frames with the same alternating Gram matrix. Here is the extension fact used to identify that orbit. In the span of any frame, choose a symplectic basis for a nondegenerate summand and a basis for its radical. Pair the radical basis with dual vectors in the orthogonal complement of the nondegenerate summand, adjust those vectors to make the pairs hyperbolic, and complete to a symplectic basis of $V$. An isometry between the original spans carries these choices to compatible bases. The resulting basis map extends it to all of $V$. Thus the extension fact also covers a degenerate restricted form.

Suppose $i$ independent images have already been selected. The conditions
\[
 \langle y_j,y_{i+1}\rangle=\langle x_j,x_{i+1}\rangle
 \quad(1\le j\le i)
\]
are $i$ independent linear equations. They have $2^{N-i}$ solutions. At most $2^i$ lie in the previous span, so there are at least $2^{N-i}-2^i$ admissible next vectors. Consequently the orbit contains at least
\begin{align*}
 \prod_{i=0}^{s-1}(2^{N-i}-2^i)
 &=2^{Ns-s(s-1)/2}\prod_{i=0}^{s-1}(1-2^{2i-N})\\
 &\ge\frac23\,2^{Ns-s(s-1)/2}
\end{align*}
frames. Indeed $s\le N/2$ makes the finite product at least
$\prod_{j\ge1}(1-4^{-j})\ge1-\sum_{j\ge1}4^{-j}=2/3$.
There are at most $2^{ts}$ ordered frames in $T$, so
\[
 \Prob(SL\subseteq T)
 \le\frac32\,2^{ts-Ns+s(s-1)/2}.
\]

The number of $s$-dimensional subspaces of $U$ satisfies
\[
 {d\brack s}_2
 =\prod_{i=0}^{s-1}\frac{2^d-2^i}{2^s-2^i}
 \le4\,2^{s(d-s)}.
\]
For the constant, $\prod_{j\ge1}(1-2^{-j})>1/4$; for example its first two factors times the elementary tail bound give
$(1/2)(3/4)(1-\sum_{j\ge3}2^{-j})=9/32>1/4$.
A union bound over the candidate subspaces gives the stated exponent and constant $6$.
\end{proof}

This is the useful transfer from kernel counting to an entire distribution: an $s$-dimensional kernel supplies a subspace, while the orbit of a whole independent frame is much larger than the orbit of one vector. The quadratic exponent in $s$ pays for the later union over every principal submatrix.

\subsection{The canonical vector in characteristic two}

Put $c=(1,\ldots,1)\in\F_2^n$ and write $B(x,y)=x^{\mathsf T}y$. Then
\[
 B(x,x)=B(x,c).
\]
Every isometry fixes $c$: for all $x$,
$B(Qx,c)=B(Qx,Qx)=B(x,x)=B(x,c)$, and nondegeneracy gives $Q^{-1}c=c$.

For $J\subseteq[n]$ let $E_J$ be its coordinate subspace, and put $j=|J|$. A vector in the kernel of $Q[J]$ is precisely a vector in
\begin{equation}\label{eq:kernel-intersection}
 K_J=E_J\cap Q^{-1}E_{J^c}.
\end{equation}
Therefore
\[
 \dim K_J\le\min(j,n-j)\le\lfloor n/2\rfloor.
\]
The empty and full sets have zero nullity, so only nonempty proper $J$ matter.

\begin{theorem}[Uniform tail for the maximum principal nullity]\label{thm:binary}
For a uniform $Q\in\mathcal O_n$ and every integer $r$ with
$1\le r\le\lfloor n/2\rfloor$,
\[
 \Prob\{\nu(Q)\ge r\}\le
 \min\{1,\;6\,2^{n-r(r-1)/2}\}.
\]
For $r>\lfloor n/2\rfloor$ the probability is zero.
\end{theorem}

\begin{proof}
We prove a bound $6\,2^{-r(r-1)/2}$ for each fixed nonempty proper $J$, then sum over at most $2^n$ choices.

\emph{Odd $n$.}
Now $B(c,c)=1$ and
\[
 \F_2^n=\langle c\rangle\perp H,\qquad H=c^\perp.
\]
The form on $H$ is alternating and nondegenerate, of dimension $N=n-1$. Restriction identifies $\mathcal O_n$ with $\Sp(H)$: every symplectic map extends by fixing $c$. A uniform $Q$ therefore induces a uniform $S\in\Sp(H)$.

Let $\rho$ be projection to $H$ along $c$. Since $c$ belongs to neither proper coordinate space, $\rho$ is injective on both $E_J$ and $E_{J^c}$. Their images have dimensions $d=j$ and $t=n-j$, so $d+t=N+1$. The projection commutes with $Q$. If $\dim K_J\ge r$, an $r$-dimensional subspace of $\rho(E_J)$ is mapped by $S$ into $\rho(E_{J^c})$. Apply Lemma~\ref{lem:symp-incidence} with $s=r\le N/2$. It gives
\[
 6\,2^{r-r(r+1)/2}=6\,2^{-r(r-1)/2}.
\]
Ignoring the remaining $c$-component imposes fewer conditions, which is legitimate for this upper bound.

\emph{Even $n$.}
Here $c\in H=c^\perp$ and the radical of $B|_H$ is exactly $\langle c\rangle$. The quotient
$\overline H=H/\langle c\rangle$ is a nondegenerate symplectic space of dimension $N=n-2$.
The natural map $\mathcal O_n\to\Sp(\overline H)$ is onto. To see this directly, choose $f=e_1$. The plane $D=\langle c,f\rangle$ has Gram matrix
$\left(\begin{smallmatrix}0&1\\1&1\end{smallmatrix}\right)$ and is nondegenerate. Its orthogonal complement is alternating and identifies with $\overline H$. Every symplectic map of that complement extends by the identity on $D$. Hence the induced $S$ is uniform.

The spaces $E_J\cap H$ and $E_{J^c}\cap H$ have dimensions $j-1$ and $n-j-1$; both inject into the quotient. Their quotient images $U,T$ satisfy $d+t=N$.
If $\dim K_J\ge r$, the subspace $K_J\cap H$ has dimension at least $r-1$. Its image in $\overline H$ remains injective and is mapped into $T$. Apply Lemma~\ref{lem:symp-incidence} with $s=r-1\le N/2$. The bound is
\[
 6\,2^{-(r-1)r/2}.
\]
When $r=1$ this step is just the trivial bound by $6$; no zero-dimensional orbit count is required.

Both parity cases give the same fixed-$J$ estimate. The union bound and the deterministic restriction following \eqref{eq:kernel-intersection} finish the proof.
\end{proof}

\subsection{Consequences and precise limits}

\begin{corollary}[No linear maximum with high probability]\label{cor:no-linear}
For fixed $0<c\le1/2$,
\[
 \Prob\{\nu(Q)\ge cn\}
 \le 6\,2^{-c^2n^2/2+(1+c/2)n}
\]
whenever the integer threshold is feasible. If $c>1/2$ the event is empty. In particular the high-probability linear lower bound proposed in~\cite{vaas} is impossible.
\end{corollary}
\begin{proof}
Use $r=\lceil cn\rceil$ in Theorem~\ref{thm:binary} and
$r(r-1)\ge cn(cn-1)$.
\end{proof}

\begin{corollary}[A square-root upper scale]
For every $a>0$ put
\[
 r_n=\left\lceil\frac{1+\sqrt{1+8(1+a)n}}2\right\rceil.
\]
Then $\Prob(\nu(Q)\ge r_n)\le6\,2^{-an}$, with probability zero if the threshold exceeds $n/2$. Consequently $\nu(Q)=O(\sqrt n)$ with exponentially high probability. More specifically, for every $\epsilon>0$, the probability that
$\nu(Q)>(\sqrt2+\epsilon)\sqrt n$ tends to zero exponentially in $n$.
\end{corollary}

The proof does not supply a matching $\Omega(\sqrt n)$ lower bound. It leaves open whether the actual scale is logarithmic, square-root, or another sublinear scale. It also does not treat independent-entry random matrices as interchangeable with orthogonal matrices: the fixed vector $c$ and the even-dimensional quotient are essential.

\subsection{Exact finite verification}

The accompanying verifier constructs orthonormal ordered bases recursively over $\F_2$, interprets each as an orthogonal matrix, and computes every principal-submatrix rank using exact bit operations. It checks dimensions $1$ through $6$, the group orders, the deterministic maximum $\lfloor n/2\rfloor$, and the stated tail bound. The finite bounds are intentionally conservative at these sizes; their purpose is to check the group convention, matrix/kernel interpretation, and edge cases. The all-$n$ result rests on the proof above. Counts and complete nullity histograms are in \path{checks/binary_orthogonal_checks.json}.

\section{Sharp finite-precision recovery along Lubin--Tate branches}
\label{sec:jets}

\subsection{The question left by the preceding report}

Berger and Sprang recover bivariate Taylor coefficients from substitutions along finitely many Lubin--Tate branches. Their Lemma~4.2 gives a finite truncation by a compactness argument, and their Proposition~4.4 uses it in coefficient extraction.\cite{bergersprang}
The earlier ProveIt packet computes the exact residue-field cutoff and characterizes the cutoff modulo powers of the uniformizer by an inverse pole. It gives the general bound
\begin{equation}\label{eq:old-jet-bound}
 t_b\le W_b+(r-1)\max_{c\ne b}(q^{v_\pi(b-c)}-1),
\end{equation}
then explicitly asks which multibranch configurations attain it, since the highest inverse terms can cancel.\cite{priorjets}
We sharpen the bound, characterize equality exactly, and give an exact all-precision formula for contacts of depths zero and one.

This is a local quantitative problem. It is distinct from the global bounded-functions conjecture motivating Berger--Sprang, for which subsequent work announces proofs. Our result neither relies on nor claims to supersede those global results.\cite{bergersprang,colovic}

\subsection{Setup and the exact interpolation threshold}

Let $K/\Q_p$ be finite, let $\OO$ be its ring of integers, and fix a uniformizer $\pi$. Normalize $v_\pi(\pi)=1$, write $q$ for the residue cardinality, and put $e=v_\pi(p)$. Fix the Lubin--Tate coordinate
\begin{equation}\label{eq:lt-coordinate}
 [\pi](X)=\pi X+X^q.
\end{equation}
All pole formulas below are proved in this coordinate; an integral conjugating change of coordinate preserves them, as explained at the end of the section.
Let $B\subset\OO$ be a finite set of $d\ge1$ distinct directions. For $r\ge1$ set $A_r=\OO/\pi^r\OO$ and
\[
 u_b(X)=[b](X)\bmod\pi^r\in XA_r[[X]],\quad
 P_b(X,Y)=\prod_{c\in B\setminus\{b\}}(Y-u_c(X)),\quad
 \Delta_b(X)=P_b(X,u_b(X)).
\]
Each branch difference is invertible in $A_r((X))$. Define
\begin{equation}\label{eq:inverse-pole}
 t_b(r)=\min\{t\ge0:X^t/\Delta_b\in A_r[[X]]\}.
\end{equation}
For $d=1$, the empty product is $1$ and $t_b(r)=0$.

\begin{proposition}[Exact componentwise recovery threshold]\label{prop:jet-threshold}
Observing $H(X,u_b(X))$ modulo $X^{n_b}$ for every $b\in B$ determines all coefficients of total degree $<d$ in every $H\in A_r[[X,Y]]$ if and only if
\[
 n_b\ge1+t_b(r)\quad\text{for every }b.
\]
The least common precision is $1+\max_b t_b(r)$. Precision $n$ retains the coefficients of $X^0,\ldots,X^{n-1}$.
\end{proposition}

\begin{proof}
Put $P=\prod_{b\in B}(Y-u_b(X))$. Division by this monic distinguished polynomial over the complete ring $A_r[[X]]$ gives a remainder $R$ of $Y$-degree $<d$. Since $P\in(X,Y)^d$, $H$ and $R$ have the same coefficients of total degree $<d$. Over $A_r((X))$, Lagrange interpolation is valid because every $u_b-u_c$ is a unit:
\[
 R(X,Y)=\sum_{b\in B}\frac{H(X,u_b(X))}{\Delta_b(X)}P_b(X,Y).
\]
If all observed residues vanish and $n_b\ge1+t_b(r)$, each quotient belongs to $XA_r[[X]]$. Since $P_b\in(X,Y)^{d-1}$, the remainder belongs to $(X,Y)^d$.

For necessity fix $b$ and set
$g_b=X^{t_b(r)}/\Delta_b\in A_r[[X]]$. Its constant coefficient is nonzero. Otherwise it would be divisible by $X$, contradicting the minimality of $t_b(r)$ (the case $t_b=0$ is immediate from the inverse identity). The series
\[
 H_b(X,Y)=g_b(X)P_b(X,Y)
\]
has nonzero coefficient $g_b(0)$ at $Y^{d-1}$, vanishes on every other branch, and evaluates to $X^{t_b(r)}$ on branch $b$. If $n_b\le t_b(r)$, all its observations vanish, so recovery fails. The argument only needs $g_b(0)\ne0$, not that it be a unit. It therefore remains valid in the presence of nilpotents in $A_r$.
\end{proof}

\subsection{Contact depths determine the inverse pole}

Fix $b$ and suppose $d\ge2$. Write
\[
 a_c=v_\pi(b-c),\qquad
 W=\sum_{c\ne b}q^{a_c},\qquad
 A=\max_{c\ne b}a_c.
\]
Define $F_a=[\pi^a](X)$, so
\begin{equation}\label{eq:lt-recurrence}
 F_0=X,\qquad F_a=\pi F_{a-1}+F_{a-1}^{\,q}.
\end{equation}

\begin{lemma}[Removing all unit parts]\label{lem:associate}
The unreduced denominator $\prod_{c\ne b}([b](X)-[c](X))$ in $\OO[[X]]$ is a unit times $\prod_{c\ne b}F_{a_c}$. Reducing modulo $\pi^r$ gives the same assertion for $\Delta_b$. Thus, for fixed $K,\pi$ and the coordinate~\eqref{eq:lt-coordinate}, $t_b(r)$ depends on the contact multiset $\{a_c:c\ne b\}$ and not on the unit parts of the differences $b-c$.
\end{lemma}
\begin{proof}
For the formal group law $\mathcal F$, the difference $\mathcal F(U,V)-U$ equals $V$ times a power series with constant term $1$. Consequently $[b]-[c]$ is a unit multiple of $[b-c]$. If $b-c=\pi^a v$ with $v\in\OO^\times$, then
$[b-c]=[v](F_a)=F_a V_a(X)$ with $V_a(0)=v$ a unit. Multiplying gives the assertion.

Multiplication by a unit power series preserves the smallest exponent of any nonzero Laurent series over $A_r$. At that exponent the coefficient is multiplied by the unit constant term and remains nonzero. This argument works even when that coefficient is nilpotent.
\end{proof}

\subsection{The leading Newton face}

\begin{lemma}\label{lem:newton-face}
For $a\ge1$ put $\lambda_a=(q-1)q^{a-1}$ and write
$F_a=\sum_jc_{a,j}X^j$. The polynomial is monic of degree $q^a$ and reduces to $X^{q^a}$ modulo $\pi$. Every nonzero lower coefficient satisfies
\[
 q^a-j\le\lambda_a v_\pi(c_{a,j}).
\]
Among lower coefficients equality occurs only at $j=q^{a-1}$, where
$c_{a,j}\equiv\pi\pmod{\pi^2}$. Thus the two terms on this leading face are exactly
\[
 X^{q^a}+\pi X^{q^{a-1}}.
\]
\end{lemma}

\begin{proof}
The case $a=1$ is immediate. For the inductive step note that
$\lambda_a=q\lambda_{a-1}$ and use~\eqref{eq:lt-recurrence}. In a nonleading product monomial from $F_{a-1}^{q}$, the total degree deficit from $q^a$ is at most $k\lambda_{a-1}$, where $k\ge1$ is the sum of the valuations of the chosen coefficients. Multinomial coefficients can only increase valuation. Hence this term is strictly above the proposed face of slope $\lambda_a$.

A monomial $\pi c_jX^j$ from the first summand has valuation $1+k$, where $k=v_\pi(c_j)$, and deficit
\[
 q^a-j=\lambda_a+(q^{a-1}-j)
 \le\lambda_a+k\lambda_{a-1}.
\]
This is strictly smaller than $(1+k)\lambda_a$ when $k>0$. When $k=0$ it is precisely $\pi X^{q^{a-1}}$. All other contributions at that exponent have valuation at least two, so they cannot cancel its residue after division by $\pi$. Adding coefficients elsewhere can only increase their valuations. This proves the assertions.
\end{proof}

\subsection{A sharp bound and a complete equality test}

\begin{theorem}[Finite-precision pole and carry criterion]\label{thm:jet-carry}
If $A=0$, then $t_b(r)=W=d-1$ for every $r$. If $A\ge1$, put
\[
 m=\#\{c\ne b:a_c=A\},\qquad \lambda=(q-1)q^{A-1}.
\]
Then
\begin{equation}\label{eq:new-jet-bound}
 \boxed{\quad t_b(r)\le W+(r-1)\lambda.\quad}
\end{equation}
Equality holds if and only if
\begin{equation}\label{eq:carry-test}
 \binom{m+r-2}{r-1}\not\equiv0\pmod p,
\end{equation}
equivalently, the addition of $m-1$ and $r-1$ in base $p$ is carry-free.
\end{theorem}

\begin{proof}
By Lemma~\ref{lem:associate}, compute with the product of the $F_{a_c}$. For $a\ge1$ write
$F_a=X^{q^a}(1+G_a)$, where $G_a$ is a Laurent polynomial all of whose coefficients are divisible by $\pi$. Modulo $\pi^r$,
\[
 (1+G_a)^{-1}=\sum_{k=0}^{r-1}(-G_a)^k.
\]
Every monomial of coefficient valuation $k$ in this inverse has exponent at least $-k\lambda_a$. At the common slope $\lambda=\lambda_A$, all nonconstant contributions from $a<A$ lie strictly above the boundary. For $a=A$ the only boundary term of $G_a$ is $\pi X^{-\lambda}$, with its coefficient interpreted modulo $\pi^2$. This is Lemma~\ref{lem:newton-face}.

The product inverse can therefore have no nonzero exponent below $-W-(r-1)\lambda$. To reach exactly that exponent with coefficient valuation below $r$, a term must use precisely the boundary contributions of the $m$ maximal-depth factors. Its coefficient modulo $\pi^r$ is consequently the coefficient supplied by
\[
 X^{-W}(1+\pi X^{-\lambda})^{-m},
\]
namely
\[
 (-1)^{r-1}\pi^{r-1}\binom{m+r-2}{r-1}.
\]
It is nonzero in $A_r$ exactly when the integer binomial coefficient is prime to $p$. This proves the inequality and equality test.

For completeness, the carry statement follows from the digit identity
\[
 v_p\binom{u+v}{u}
 =\frac{s_p(u)+s_p(v)-s_p(u+v)}{p-1},
\]
where $s_p$ is the sum of base-$p$ digits. This identity follows by applying
$v_p(n!)=\sum_{h\ge1}\lfloor n/p^h\rfloor=(n-s_p(n))/(p-1)$ to the three factorials. Its value is zero exactly when the addition has no carry. If $A=0$, every $F_{a_c}=X$ and the conclusion is immediate.
\end{proof}

\paragraph{Exactly when the previous bound can be sharp.}
The preceding estimate~\eqref{eq:old-jet-bound} is $W+(r-1)(q^A-1)$. The new theorem gives a full componentwise classification of its attainment:

\begin{center}\small
\begin{tabular}{ll}
\toprule
Case & Attainment of the preceding bound\\
\midrule
$r=1$ & Always.\\
$A=0$ & Always, at every $r$.\\
$r>1$, $A=1$ & Exactly when $(m-1)+(r-1)$ has no base-$p$ carry.\\
$r>1$, $A\ge2$ & Never: $q^A-q^{A-1}<q^A-1$.\\
\bottomrule
\end{tabular}
\end{center}

Thus the question left by the prior report is answered for each branch. The sharpened bound itself has the exact carry criterion at every depth. When that criterion fails and $A\ge2$, lower Newton faces can determine the precise smaller pole; we do not claim a closed formula for all those remaining values.

\begin{corollary}[Two branches at arbitrary contact depth]\label{cor:two-branches}
For exactly two branches with $a=v_\pi(b-c)\ge1$,
\[
 t_b(r)=q^a+(r-1)(q-1)q^{a-1},
\]
and both branches have the same least precision $1+t_b(r)$.
\end{corollary}
\begin{proof}
There is one maximal-depth factor, so $m=1$ and the binomial coefficient is $1$ at every $r$.
\end{proof}

\subsection{An exact formula when all contacts have depth at most one}

\begin{theorem}[All-precision depth-one formula]\label{thm:depth-one}
Suppose every $a_c\in\{0,1\}$ and let $m\ge1$ count the contacts of depth one. Set
\[
 K_m(r)=\max\left\{0\le k<r:
 k+e\,v_p\binom{m+k-1}{k}<r\right\}.
\]
Then
\begin{equation}\label{eq:exact-depth-one}
 \boxed{\quad t_b(r)=W+(q-1)K_m(r).\quad}
\end{equation}
The candidate set is nonempty, since it contains $k=0$.
\end{theorem}
\begin{proof}
The unit-free denominator is exactly
\[
 X^W(1+\pi X^{-(q-1)})^m.
\]
Its reciprocal is
\[
 X^{-W}\sum_{k=0}^{r-1}(-1)^k
 \binom{m+k-1}{k}\pi^kX^{-k(q-1)}
 \pmod{\pi^r}.
\]
Distinct $k$ give distinct exponents. The coefficient at index $k$ has valuation
$k+e\,v_p\binom{m+k-1}{k}$, so it survives exactly under the condition in the maximum. No cancellation between different $k$ is possible.
\end{proof}

\paragraph{A complete three-branch example.}
Take $K=\Q_3$, $\pi=3$, $q=3$, and $B=\{0,3,6\}$. Every pair of directions has contact depth one. At each branch $W=6$ and $m=2$. For $r=3$, the indices $k=0,1$ survive, whereas the coefficient at $k=2$ is
$3^2\binom32=27$, which vanishes modulo $27$. Thus
\[
 t_b(3)=8,\qquad \text{least common precision}=9.
\]
The prior upper bound gives precision $11$. Observations through degree eight suffice; through degree seven they do not, by the explicit interpolation witness of Proposition~\ref{prop:jet-threshold}.

\begin{center}\small
\begin{tabular}{c|rrrrrrrrr}
\toprule
$r$ &1&2&3&4&5&6&7&8&9\\
\midrule
$K_2(r)$ &0&1&1&3&4&4&6&7&7\\
Exact precision $1+t_b(r)$ &7&9&9&13&15&15&19&21&21\\
Prior upper precision &7&9&11&13&15&17&19&21&23\\
\bottomrule
\end{tabular}
\end{center}

The plateau at a multiple of three is an arithmetic cancellation, not a rounding effect. For comparison, two branches with contact depth two at $q=3,r=3$ have precision
$1+9+2\cdot6=22$, improving the earlier upper precision $1+9+2\cdot8=26$ even without multibranch cancellation.

\subsection{A bounded and eventually periodic precision defect}

The exact formula also controls its dependence on the requested precision. This gives an additional computable consequence without a new global arithmetic input.

\begin{corollary}\label{cor:period}
In Theorem~\ref{thm:depth-one}, let $P$ be the least power of $p$ with $P\ge m$. For every $r\ge1$,
\[
 \max(0,r-P)\le K_m(r)\le r-1.
\]
For $m\ge2$ set $d=m-1$,
\[
 L=\left\lfloor\frac{P-1}{e}\right\rfloor,
 \quad T=L+v_p(d!)+1,\quad Q=p^T.
\]
Then, for $r\ge P$,
\[
 K_m(r+Q)=K_m(r)+Q.
\]
Thus $t_b(r+Q)=t_b(r)+(q-1)Q$. For $m=1$ the exact formula is $K_1(r)=r-1$ from the start.
\end{corollary}

\begin{proof}
Every $k$ divisible by $P$ has no carry when added to $m-1<P$, so its binomial coefficient is a $p$-adic unit. The largest such $k\le r-1$ proves the first assertion.
For $r\ge P$, write a potential maximum as $k=r-1-\delta$ with $0\le\delta<P$. It is admissible exactly when
\[
 e\,v_p\binom{k+d}{d}\le\delta.
\]
All thresholds for the binomial valuation are at most $L$. Since
$\binom{k+d}{d}=\prod_{i=1}^d(k+i)/d!$, these comparisons are determined by $k\bmod p^T$. If a numerator valuation reaches $T$, it already exceeds every relevant threshold $L+v_p(d!)$; below that threshold each factor's valuation is determined by its residue modulo $p^T$. Replacing $r$ by $r+Q$ leaves the admissible $\delta$ unchanged. Its smallest value, which gives the largest $k$, is therefore unchanged.
\end{proof}

The period displayed here is a safe explicit period, not a claim of optimality. The finite maximum~\eqref{eq:exact-depth-one} is usually the simpler algorithm.

\subsection{Coordinate changes, verification, and remaining limits}

Suppose an integral invertible coordinate change $h(X)=uX+\cdots$, $u\in\OO^\times$, conjugates the formal modules. Then
\[
 h(U)-h(V)=(U-V)H(U,V),\qquad H(0,0)=u.
\]
Every transformed branch difference is therefore the old difference, after the invertible substitution in $X$, times a unit. Invertible substitution preserves the smallest Laurent exponent, including over $A_r$. Consequently the pole formulas persist under such an integral conjugacy. The total-degree filtration of $A_r[[X,Y]]$ is preserved as well. This is the exact coordinate-invariance assertion used; no arbitrary nonintegral change of parameter is allowed.

The verifier \path{code/check_jets.py} independently builds $F_a$ by recursion and inverts the normalized products by finite geometric series in sparse Laurent polynomials. It checks the full inverse-product identity, the improved upper bound, its equality criterion, strictness relative to the older bound, and the depth-one formula. All 210 exact cases passed for $(p,q)=(2,2),(2,4),(3,3),(3,9)$ at precisions through $r=6$, with the contact profiles recorded in \path{checks/jet_checks.json}. These examples use unramified coefficient arithmetic; the proof above, including the factor $e$, handles ramified fields. No finite experiment substitutes for that argument.

The new conclusion is a precise answer to a question in the supplied earlier analysis and a quantitative refinement of a published interpolation method. We do not claim that the local formulas themselves were a named longstanding conjecture, nor that a broad literature search establishes their first publication here.

\section{Exact tests for a block of a split orthogonal matrix}
\label{sec:orthogonal}

The finite Fourier argument in
\path{grh/src/New/Spanning/SecondMoment/PositiveType.lean}
is independent of the final zero-free theorem. Its estimate
\texttt{re\_sum\_addChar\_dotProduct\_mul\_le} bounds an additive character
coefficient through the masses of two stabilizers; its final theorem
\texttt{card\_mul\_sum\_det\_eq\_zero\_le} applies this to a singular
matrix compression. Earlier ProveIt reports already passed from vectors
to frames, obtained rectangular matrix distribution bounds, and computed
the precise obstruction for isotropic blocks of symplectic matrices.
We do not count those results again. Here is a different exact boundary
calculation for split orthogonal matrices. It identifies an optimal
statistical test for the \emph{entire block}, not only its rank.

This result is a complete human proof. The group counts it uses are
classical; for related counts of isotropic subspaces see Li--Zhang,
Lemma~3.2.2, and for finite-field rank laws see Fulman--Goldstein.\cite{lizhang,fulman}
The exact block comparison below is new to the inspected report corpus.
No claim of worldwide priority is made.

\subsection{Statement}

Let $q$ be an odd prime power, $m\ge1$, and put
\[
 J=\begin{pmatrix}0&I_m\\I_m&0\end{pmatrix},\qquad
 G=\{g\in\operatorname{GL}_{2m}(\mathbb F_q):g^{\mathsf T}Jg=J\}.
\]
This is the split orthogonal group $O^+(2m,q)$.
Its determinant-one subgroup is denoted by $G^+=SO^+(2m,q)$; this
notation does not mean the possibly smaller commutator subgroup.
Write
\[
 g=\begin{pmatrix}A&B\\C&D\end{pmatrix}
\]
in this hyperbolic basis, and let $U_m$ be a uniform $m\times m$ matrix
over $\mathbb F_q$. Thus the entries of $U_m$ are independent and uniform.
For laws $P,Q$ on the finite matrix space,
$d_{\mathrm{TV}}(P,Q)=\frac12\sum_M|P(M)-Q(M)|$.

\begin{theorem}[Optimal parity and full-rank tests]
\label{thm:orthogonal-optimal-test}
If $g$ is uniform in $G^+$, then $\operatorname{rank}C$ is even and
\begin{equation}
 d_{\mathrm{TV}}(C,U_m)
   =\Pr\bigl(\operatorname{rank}U_m\text{ is odd}\bigr).
 \label{eq:SO-TV}
\end{equation}
Consequently the test ``the block has odd rank'' achieves the full
total variation distance; no test using all its entries improves this
distinguishing advantage.

If instead $g$ is uniform in the full group $G$, then
\begin{equation}
 d_{\mathrm{TV}}(C,U_m)
  =\prod_{j=1}^{m}(1-q^{-j})
    -\frac{1}{2\prod_{j=1}^{m-1}(1+q^{-j})}.
 \label{eq:O-TV}
\end{equation}
Here the full-rank test achieves the full distance. Empty products equal
one, so the formula also includes $m=1$.
\end{theorem}

For clarity, ``achieves the full distance'' uses
$d_{\mathrm{TV}}(P,Q)=\max_E|P(E)-Q(E)|$.
With equally likely hypotheses, the corresponding best success
probability is $(1+d_{\mathrm{TV}})/2$.

\subsection{Proof by counting isotropic subspaces}

Let $W=\mathbb F_q^m\oplus0$ and $L=gW$.
The space $L$ is maximal totally isotropic for $J$ and
\[
 \operatorname{rank}C=m-\dim(L\cap W).
\]
Write
\[
 {m\brack r}_q
 =\prod_{i=0}^{r-1}\frac{q^m-q^i}{q^r-q^i},\quad
 F_q(m,r)=\prod_{i=0}^{r-1}(q^m-q^i),\quad
 T_m=2\prod_{i=1}^{m-1}(1+q^i).
\]

\paragraph{The determinant and rank parity.}
Every element of $G$ stabilizing $W$ has block form
\(
 \left(\begin{smallmatrix}P&*\\0&P^{-\mathsf T}\end{smallmatrix}\right)
\), hence determinant one.
Suppose $\dim(L\cap W)=m-r$. A change of basis in $W$, extended by the
inverse transpose on its complementary isotropic subspace, puts the
intersection in standard position. In its orthogonal quotient, $L$ is
transverse to the remaining part of $W$, hence is the graph of a
skew-symmetric matrix. An upper unipotent isometry with a skew-symmetric
upper-right block removes that graph. Both changes stabilize $W$ and
have determinant one. The resulting space is
\[
 \langle e_1,\ldots,e_{m-r},f_{m-r+1},\ldots,f_m\rangle.
\]
It is the image of $W$ under the isometry which swaps $e_i$ with $f_i$
in precisely $r$ hyperbolic pairs. That isometry has determinant
$(-1)^r$. Multiplying by it leaves an element stabilizing $W$, proving
\begin{equation}
 \det g=(-1)^{\operatorname{rank}C}.
 \label{eq:orthogonal-parity}
\end{equation}
This construction also proves transitivity on maximal isotropic
subspaces. Since their stabilizer has determinant one, the two determinant
cosets of $G$ give two equally sized families, selected by the parity
of $r$.

\paragraph{The exact rank counts.}
Choose $D=L\cap W$ of dimension $m-r$ in ${m\brack r}_q$ ways.
The quotient $D^\perp/D$ is a hyperbolic space of dimension $2r$.
The possible $L/D$ are precisely its maximal isotropic complements to
$W/D$. Writing each such complement as a graph gives a skew-symmetric
$r\times r$ matrix. There are $q^{r(r-1)/2}$ such matrices because
$q$ is odd. Thus the number of possible $L$ at block rank $r$ is
\begin{equation}
 N_m(r)={m\brack r}_q q^{r(r-1)/2}.
 \label{eq:orthogonal-subspace-count}
\end{equation}
The elementary $q$-binomial identity
\[
 \sum_{r=0}^{m}{m\brack r}_q q^{r(r-1)/2}
   =\prod_{i=0}^{m-1}(1+q^i)=T_m
\]
follows by induction from the Gaussian-binomial recursion.
Uniform group elements induce the uniform law on each orbit. Therefore
the full orthogonal rank law and the special orthogonal rank law are
\begin{align}
 o_m(r)&=N_m(r)/T_m,\label{eq:O-rank-law}\\
 s_m(r)&=\begin{cases}2o_m(r),&r\text{ even},\\0,&r\text{ odd}.
                 \end{cases}\label{eq:SO-rank-law}
\end{align}
For an independent-entry matrix, choose its $r$-dimensional image and
then a surjection onto it. This gives
\begin{equation}
 u_m(r)=\Pr(\operatorname{rank}U_m=r)
       ={m\brack r}_qF_q(m,r)q^{-m^2}.
 \label{eq:iid-rank-law}
\end{equation}

\paragraph{Why rank suffices for all matrix tests.}
For each $P\in\operatorname{GL}_m(\mathbb F_q)$,
the matrix $\operatorname{diag}(P,P^{-\mathsf T})$ lies in $G^+$.
Left and right multiplication by these matrices transform $C$ into
$P^{-\mathsf T}CQ$ for independently arbitrary invertible $P,Q$.
Uniform measure on either group is invariant under these operations.
Since left and right invertible transformations act transitively on
matrices of a given rank, both block laws are uniform conditional on
rank. The independent-entry law has the same property. Thus the total
variation distance of the full matrices equals that of their rank laws.

\paragraph{Comparing the mass at each rank.}
Set $R_m(r)=2o_m(r)/u_m(r)$, also for odd $r$. Cancellation in the exact
formulas gives
\[
 R_m(m)=\frac{1}{(1-q^{-m})
                      \prod_{j=1}^{m-1}(1-q^{-2j})}>1,
 \qquad
 \frac{R_m(r-1)}{R_m(r)}=q^{m-r+1}-1.
\]
Because $q\ge3$, each displayed ratio is at least two.
Hence $s_m(r)>u_m(r)$ at every even rank, while $s_m(r)=0$ at every
odd rank. All the deficit lies at the odd ranks, proving
\eqref{eq:SO-TV}.

For the full orthogonal group, the likelihood ratio is $R_m(r)/2$.
At rank $m$ it is less than one: writing $Q=q^{-1}\le1/3$, its
denominator before the factor two satisfies
\[
 (1-Q^m)\prod_{j=1}^{m-1}(1-Q^{2j})
 \ge (1-Q)\left(1-\sum_{j=1}^{\infty}Q^{2j}\right)
 \ge \frac{7}{12}>\frac12.
\]
It is greater than $1/2$, because that same product is less than one.
At rank $m-1$ the ratio is therefore greater than
$(q-1)/2\ge1$, and it increases as rank decreases.
Thus only full rank has deficient mass in $o_m$ compared with $u_m$.
Their full-rank probabilities are respectively
\[
 o_m(m)=\frac{1}{2\prod_{j=1}^{m-1}(1+q^{-j})},\qquad
 u_m(m)=\prod_{j=1}^{m}(1-q^{-j}).
\]
Their difference proves~\eqref{eq:O-TV}.\hfill$\square$

\subsection{Finite examples and the limiting obstruction}

The supporting program \path{code/check_orthogonal_blocks.py} uses exact rational
arithmetic. It independently enumerates all maximal isotropic subspaces
and all square matrices for $q=3$, $m=1,2,3$. At $m=3$ it examines all
$33{,}880$ three-dimensional subspaces of $\mathbb F_3^6$ in unique
reduced-row-echelon form, of which $80$ are isotropic. It also independently
enumerates all $2$ and $576$ elements of $SO^+(2m,3)$ for $m=1,2$ from
elementary generators, checks the preserved form, and checks every
individual block probability. No group enumeration at $m=3$ is claimed.

\begin{center}
\begin{tabular}{c|c|c|l}
 $m$ & $d_{\mathrm{TV}}(C_{SO},U_m)$ &
       $d_{\mathrm{TV}}(C_O,U_m)$ & nonzero $SO$ rank masses\\\hline
 $1$ & $2/3$ & $1/6$ & $r=0:\ 1$\\
 $2$ & $32/81$ & $47/216$ & $r=0:\ 1/4,\quad r=2:\ 3/4$\\
 $3$ & $11570/19683$ & $13597/58320$ &
       $r=0:\ 1/40,\quad r=2:\ 39/40$
\end{tabular}
\end{center}

For the full orthogonal group, fixed $q$ and $m\to\infty$ in
\eqref{eq:O-TV} give the positive limiting distance
\[
 \prod_{j=1}^{\infty}(1-q^{-j})
 -\frac{1}{2\prod_{j=1}^{\infty}(1+q^{-j})}.
\]
The products converge to positive limits, and the proof above shows
that the difference is positive for every odd $q$.
Even when $q\to\infty$, this expression tends to $1/2$.
The matrix entries therefore do not become independent at this maximal
isotropic boundary. The determinant component leaves a precise parity
signature, and averaging the two components leaves a precise full-rank
signature. This is a method transfer in finite geometry and probability,
with no use of RH or GRH as a premise.


\section{What this establishes, and where the larger conjectures still resist}
\label{sec:assessment}

\subsection{Claim ledger}

\begin{longtable}{P{.29\linewidth}P{.29\linewidth}P{.30\linewidth}}
\toprule
Claim & Status in this report & Boundary of the claim\\
\midrule\endhead
Binary principal-nullity tail & Complete all-dimension proof; explicit upper tail and negative answer to the linear-growth proposal & No matching lower bound or exact typical order.\\[5pt]
Lubin--Tate finite-precision bound & Complete proof and exact carry criterion for attaining the sharpened bound & For maximal contact depth $\ge2$ when the leading coefficient cancels, no general closed formula for the next pole is claimed.\\[5pt]
Depth-zero/one branch contacts & Exact formula at every coefficient precision; explicit eventual affine period & The displayed period is not asserted minimal.\\[5pt]
Split orthogonal half-block law & Exact total-variation formulas for determinant-one and full groups, for every odd $q$ and every block size & This is a boundary calculation in the specified hyperbolic basis, not a theorem for arbitrary real orthogonal blocks.\\[5pt]
Novelty relative to prior reports & The closest source files and 52-archive inventory were compared; specific duplication was excluded & Inventory and targeted reading are not a proof audit of the whole corpus or a worldwide priority determination.\\[5pt]
Formal verification & Source declarations inspected; new mathematical proofs and exact Python checks supplied & No fresh Lean build of upstream or of these new theorems.\\[5pt]
General automorphic or nonabelian RH & Precise transfer requirements identified & No new case proved here.\\
\bottomrule
\end{longtable}

The first theorem is an explicit resolution of a part of a publicly posted question. The second completes a sharpness problem deliberately left open in the supplied research corpus. Those are concrete outcomes beyond listing standard conditional consequences of RH. Their mathematical ingredients, such as Witt extension, Gaussian-binomial counts, Lagrange interpolation, and factorial valuations, are classical and retain that status.

\subsection{The strongest independent arithmetic theorem inside the source}

A particularly worthwhile extraction target is
\src{src/New/Congruences/Selmer/GeneralTheorem.lean}{\code{New.Selmer.exists\_anticyclotomic\_twist}}.
For a fixed coefficient prime, finitely many compatible Selmer data and Stickelberger systems, and specified torsion-level hypotheses, it constructs anticyclotomic characters whose first several powers simultaneously have bounded normalized Selmer lengths, controlled valuations of the group-ring systems, and small auxiliary ramification costs. This is a substantive arithmetic result inside the proof, and it is not an ordinary consequence of RH.\cite{selmer}

The quantifiers matter. The character order, coefficient extension, and auxiliary primes can vary with the requested tolerance. The cost constants include data-dependent local and global terms; they are not automatically zero. The theorem does not prescribe a fixed character order, establish positive density in a conductor family, or supply one character simultaneously good at every coefficient prime. Its later use in a regular family intentionally permits place-dependent choices compatible with a height bound.

Nor is ``simultaneous nonvanishing for a finite family'' by itself a new target. Kriz--Nordentoft's horizontal $p$-adic $L$-functions already address simultaneous nonvanishing and Diophantine stability, while Mazur--Rubin establish broad stability results in prime-power extensions. A publishable extension would need a specified new family, an improved quantitative estimate, or genuinely stronger order/density constraints.\cite{kriznordentoft,mazurrubin}

\begin{proposition}[A normalized-length warning with an exact example]\label{prop:length-warning}
For any prime $\ell$, there are nonzero modules $M_N$ over coefficient DVRs $O_N$ such that their lengths normalized by $\operatorname{length}_{O_N}(O_N/\ell)=1$ tend to zero.
\end{proposition}
\begin{proof}
Take $O_N=\Z_\ell[\zeta_{\ell^N}]$, with uniformizer $\varpi_N=\zeta_{\ell^N}-1$, and put $M_N=O_N/(\varpi_N)$. The ramification index is
$e_N=\ell^{N-1}(\ell-1)$. The ordinary lengths of $M_N$ and $O_N/\ell$ are $1$ and $e_N$, respectively. Thus the normalized length is $1/e_N\to0$, while every $M_N$ is nonzero. For $\ell=2$ one can start at $N=2$.
\end{proof}

This refutes only the inference from shrinking normalized bounds over increasingly ramified coefficient rings to vanishing. It does not assert that any actual Selmer group produced by the source is nonzero. Over a fixed coefficient DVR, an ordinary nonnegative integer length smaller than one would of course vanish. A transfer to BSD, Goldfeld, or a fixed-order twist question must first supply the missing comparison and uniformity.

\subsection{A quantitative test for a proposed global extension}

The essential screening question is whether the new arithmetic construction supplies both bounds on the same object, with a loss small enough after division by the amplification parameter. One useful algebraic model is the following; it is an abstraction, not a new claim of historical novelty.

\begin{proposition}[Spanning fraction and error budget]\label{prop:budget}
Let $D,n,m,\eta>0$, with $\eta D\le n\le D$. Let $U,L,I\ge0$, and suppose $A+U\ge0$ and real quantities $V,\delta,R$ satisfy
\[
 D(A-L)\le V+2\delta,\qquad
 V\le n(A-mR+U),\qquad \delta\le DI.
\]
Then
\[
 R\le\frac{U+L+2I}{\eta m}.
\]
\end{proposition}
\begin{proof}
Combining the three inequalities gives
\[
 nmR\le(n-D)(A+U)+D(U+L+2I)
 \le D(U+L+2I).
\]
Divide by $nm>0$ and use $n\ge\eta D$ and nonnegativity of the numerator.
\end{proof}

In an actual theorem, the cost of completing the independent family must be included in $U,L,I$ or in the comparison itself. The near-full spanning fraction in the source cannot be replaced casually by a fixed positive fraction without paying that cost. If an extension gives $m\ge x^\alpha$, $\eta\ge x^{-\lambda}$, and a normalized error budget $x^{u+o(1)}$, this model yields at most
$R(x)\le x^{-\alpha+\lambda+u+o(1)}$. Square-root strength requires $\alpha-\lambda-u\ge1/2$. This is a useful way to reject an attractive but quantitatively insufficient construction before completing a large formal development.

\begin{center}\small
\begin{tabular}{P{.21\linewidth}P{.33\linewidth}P{.36\linewidth}}
\toprule
Target & Reusable source mechanism & Decisive missing result\\
\midrule
Hecke RH over a new base field & Theta norms, twisting, determinant framework & Prime-ideal divisibility and period cancellation with the correct errors.\\[5pt]
RH for a fixed non-CM $\GL_2$ form & Norm identities and congruence control & An arithmetic detector retaining that form's zero information and a full matching comparison.\\[5pt]
Nonabelian Artin RH & Representation separation and one-sided transforms & A Frobenius-sensitive arithmetic estimate; cancellation and possible poles must also be controlled.\\[5pt]
Fixed-order Selmer/nonvanishing conjectures & Simultaneous anticyclotomic specialization & Fixed order, fixed coefficients or a uniform descent theorem, and conductor/density control as appropriate.\\[5pt]
Simplicity and pair correlation & The source controls real parts of zeros & New information about multiplicities or correlations; zero exclusion alone does not provide it.\\
\bottomrule
\end{tabular}
\end{center}

For larger zero-location conjectures, the most promising next step is an arithmetic construction for one fully specified target, with its error budget written in advance. More equivalent prime-sum criteria do not resolve that task. The independent finite and local theorems in this report demonstrate a second productive route: isolate a precise internal method, adapt its hypotheses, and solve a concrete problem for which all required estimates can be proved.

\subsection{Reproducibility and formalization}

The ZIP contains the complete \LaTeX{} source, compiled PDF, four new Python checkers, their machine-readable results, a source manifest, and a formalization plan. All checkers use only Python's standard library. The main verifier examined 23,817 binary orthogonal matrices, 1,498,426 principal submatrices, and 238,557 independent symplectic frames. An independently written binary enumerator reproduced the group orders and nullity histograms. The local checker performed 210 exact Laurent inversions. The odd-characteristic block checker enumerated all 33,880 three-dimensional subspaces of $\F_3^6$ and all 19,683 matrices of size three, with direct group checks at sizes one and two.

The proofs establish the infinite families. The computations check the definitions, arithmetic, and small examples; their scope is recorded explicitly in each JSON file. A source manifest records upstream paths and hashes, and the prior-archive manifest records original download identities. The previous archives themselves are not duplicated inside this deliverable.

The formalization plan separates the finite symplectic incidence lemma, the characteristic-two quotient, the coefficient valuation face, and the interpolation theorem over a nonreduced ring. It names mathematical interfaces rather than presenting uncompiled sketches as completed Lean code. No new theorem in this report is represented as accepted by Lean's kernel.

\begingroup
\setlength{\parskip}{0pt}
\begin{thebibliography}{99}\small\raggedright
\setlength{\itemsep}{3pt}
\setlength{\parsep}{0pt}
\bibitem{repo} The \texttt{lsernau/grh-lean} contributors.
\emph{GRH for Dirichlet $L$-functions, in Lean 4}.
\href{https://github.com/lsernau/grh-lean/tree/72fead102f663d5c7a407964a13dc0cf9d977207}{Repository at commit 72fead102f663d5c7a407964a13dc0cf9d977207}, 25 November 2026. The repository attributes the underlying method to Levent Alp\"oge with Claude. See its README for reported validation and its Apache-2.0 license.

\bibitem{statements} Upstream \src{src/Statements.lean}{\texttt{Statements.lean}} and \src{src/Main.lean}{\texttt{Main.lean}}, at the preceding commit. These determine the exported theorem scope.

\bibitem{guide} Upstream \src{GUIDE.md}{\emph{Reading guide}}, especially the three main estimates, glossary, parameter sizes, and conventions.

\bibitem{mainineq} Upstream \src{src/New/Reduction/MainInequality.lean}{\texttt{New/Reduction/MainInequality.lean}}, \src{src/New/Reduction/ShortAndEnough.lean}{\texttt{ShortAndEnough.lean}}, and \src{src/New/Determinant/Twist/LongForms.lean}{\texttt{LongForms.lean}}.

\bibitem{criterion} Upstream \src{src/New/Reduction/OneSidedBound.lean}{\texttt{OneSidedBound.lean}} and \src{src/New/Reduction/Zeros/Criterion.lean}{\texttt{Zeros/Criterion.lean}}.

\bibitem{positivetype} Upstream \src{src/New/Spanning/SecondMoment/PositiveType.lean}{\texttt{PositiveType.lean}}, especially \texttt{card\_mul\_sum\_det\_eq\_zero\_le}.

\bibitem{onecurve} Upstream \src{src/New/Derivatives/Sharpness/OneCurve.lean}{\texttt{OneCurve.lean}}, especially \texttt{exists\_monomial\_max\_norm\_coeff}.

\bibitem{selmer} Upstream \src{src/New/Congruences/Selmer/GeneralTheorem.lean}{\texttt{GeneralTheorem.lean}}, \texttt{New.Selmer.exists\_anticyclotomic\_twist}.

\bibitem{priorjets} ProveIt incoming corpus.
\prior{grh_transfers_cycle_signings_and_lubin_tate.zip}{\emph{Optimal signings of square cycles and sharp Lubin--Tate jet recovery}}, 25 November 2026; particularly \texttt{lubin\_tate.tex}, the finite-precision proposition, and the third concluding research question. Blob identity and SHA-256 appear in the accompanying prior manifest.

\bibitem{priorrank} ProveIt incoming corpus.
\prior{grh_beyond_dirichlet_2026-11-25.zip}{\texttt{grh\_beyond\_dirichlet\_2026-11-25.zip}}, especially \texttt{sections/rank\_moments.tex}.

\bibitem{priorsymplectic} ProveIt incoming corpus.
\prior{grh_sharp_torsion_and_symplectic_blocks.zip}{\texttt{grh\_sharp\_torsion\_and\_symplectic\_blocks.zip}}, especially \texttt{symplectic.tex}. All prior citations use snapshot \texttt{9816644452bb42827167d6110e6d0322315a02b9}.

\bibitem{vaas} Vaas.
\emph{Principal submatrices of random binary orthogonal matrices}.
\href{https://math.stackexchange.com/questions/5119604/principal-submatrices-of-random-binary-orthogonal-matrices}{Mathematics Stack Exchange question 5119604}. Retrieved 26 November 2026. Cited as the primary statement of the question, not as authority for our proof.

\bibitem{bergersprang} L. Berger and J. Sprang.
\emph{Integer-valued polynomials and $p$-adic Fourier theory}.
\href{https://arxiv.org/abs/2502.18053v2}{arXiv:2502.18053v2} (2025), Lemma 4.2 and Proposition 4.4.
\href{https://perso.ens-lyon.fr/laurent.berger/articles/POLINT.pdf}{Author PDF}.

\bibitem{bhargava} M. Bhargava.
$P$-orderings and polynomial functions on arbitrary subsets of Dedekind rings.
\emph{J. reine angew. Math.} 490 (1997), 101--127.
\href{https://doi.org/10.1515/crll.1997.490.101}{DOI: 10.1515/crll.1997.490.101}. Foundational context for effective nonarchimedean interpolation; no novelty for that general subject is asserted here.

\bibitem{colovic} U. A. Colovi\'c.
\emph{On the ring of bounded functions on the character variety}.
\href{https://arxiv.org/abs/2611.25162}{arXiv:2611.25162}, submitted 18 November 2026. Cited for its separately announced global theorem, not as a premise in our local proof.

\bibitem{lizhang} C. Li and W. Zhang.
\emph{On the arithmetic Siegel--Weil formula for GSpin Shimura varieties}.
\href{https://www.math.columbia.edu/~chaoli/KRO.pdf}{Author manuscript, 2 March 2022}, Lemma 3.2.2 for classical isotropic-subspace counts.

\bibitem{fulman} J. Fulman and L. Goldstein.
Stein's method and the rank distribution of random matrices over finite fields.
\emph{Annals of Probability} 43 (2015), 1274--1314.
\href{https://arxiv.org/abs/1211.0504}{arXiv:1211.0504};
\href{https://doi.org/10.1214/13-AOP889}{DOI: 10.1214/13-AOP889}.

\bibitem{kriznordentoft} D. Kriz and A. C. Nordentoft.
\emph{Horizontal $p$-adic $L$-functions}.
\href{https://arxiv.org/abs/2310.20678}{arXiv:2310.20678}, version 3, 16 November 2025.

\bibitem{mazurrubin} B. Mazur and K. Rubin, with an appendix by M. Larsen.
\emph{Diophantine stability}.
\href{https://arxiv.org/abs/1503.04642}{arXiv:1503.04642}, version 3, 3 July 2017.
\end{thebibliography}
\endgroup

\end{document}
